% id: 41-08
% book: 베이직쎈 확률과 통계 · 03 확률의 뜻과 활용
% section: 유형 023 배반사건
% figure: none
% answer: 
% answer_source: 
% variant_level: 0 (원본 전사)
% 이 파일은 build.mjs 가 items.json 에서 만든다. 고칠 때는 items.json 을 고친다.
\smash{\raisebox{1.5mm}{\dmkichul{베이직쎈 확률과 통계 41쪽 08번}}}
1부터 20까지의 자연수가 각각 하나씩 적힌 20장의 카드 중 임의로 한 장의 카드를 뽑는 시행에서 세 사건 $A$, $B$, $C$가 각각 다음과 같을 때, 서로 배반사건인 것만을 보기에서 있는 대로 고른 것은?\exprbox{\begin{minipage}{0.96\linewidth}\raggedright $A$: 5의 배수가 적힌 카드를 뽑는 사건\\ $B$: 9의 배수가 적힌 카드를 뽑는 사건\\ $C$: 15의 약수가 적힌 카드를 뽑는 사건\end{minipage}}\bogi{ㄱ. $A$와 $B$ \hfill ㄴ. $A$와 $C$ \hfill ㄷ. $B$와 $C$}
\dmchoicesauto{}{ㄱ}{ㄷ}{ㄱ, ㄴ}{ㄱ, ㄷ}{ㄴ, ㄷ}
